\documentclass{KPS_assignment}
% ---------For Persian------------
%\newcommand*{\name}{کیوان جمالی}
%\newcommand*{\id}{403209743}
%\newcommand*{\professor}{دکتر شفاهی}
%\newcommand*{\cdate}{2024/12/26}
%\newcommand*{\course}{درس (شماره)}
%\newcommand*{\assignment}{تمرین شماره1}
% ----------For English------------
\newcommand*{\name}{Keivan Jamali}
\newcommand*{\id}{403209743}
\newcommand*{\professor}{Prof. Amini}
\newcommand*{\cdate}{1404/(7, 10)}
\newcommand*{\course}{Advanced Traffic (20--551)}
\newcommand*{\assignment}{Class Project}
% ----------START------------------
\begin{document}
\maketitle

\section*{Project Description: Network Traffic Simulation}

\textbf{Goal:}  
Students will design, implement, and analyze a traffic network simulation using Python. The simulation represents vehicle movement in a network of nodes, edges, blocks, and intersections with traffic lights. After completing the Python-based simulation, the same network will be modeled and simulated inside SUMO to compare results.

\vspace{0.3cm}
\subsection*{Libraries Allowed}
Students may use any of the following (or similar) Python packages:
\begin{itemize}
    \item \texttt{pandas} for data reading and handling
    \item \texttt{matplotlib} for plots
    \item \texttt{networkx} for graph/network structure
    \item \texttt{simpy} for event simulation
    \item \texttt{sympy} for formulas or any other library
\end{itemize}

\section*{Part 1: Network Structure}

The input consists of:
\begin{itemize}
    \item \textbf{Node file} – list of nodes with IDs and optional coordinates.
    \item \textbf{Edge file} – for each edge: start node, end node, length, number of lanes (optional).
    \item \textbf{Demand file} – each row indicates:
    \begin{itemize}
        \item vehicle ID,
        \item start time,
        \item origin node,
        \item destination node.
    \end{itemize}
\end{itemize}

\vspace{0.2cm}
\textbf{Network configuration:}
\begin{itemize}
    \item Each edge is divided into blocks of \(100\) meters.
    \item Each block has capacity of about \(20\) vehicles.
    \item Vehicles may move one block every \(5\) simulation steps.
    \item There may be more than one lane per edge (optional part).
\end{itemize}

\section*{Part 2: Vehicle Simulation Rules}

\begin{itemize}
    \item Simulation runs in discrete time steps (\(1\) second per step).
    \item Vehicles are not required to enter only at times \(5, 10, 15...\); they may appear at any time (e.g., \(t=7\)).
    \item If a block is full, vehicles must wait and accumulate delay time.
    \item Each intersection has a traffic signal:
    \begin{itemize}
        \item fixed-time configuration,
        \item or adaptive, (optional)
        \item or any custom logic.
    \end{itemize}
    \item Vehicles must obey signal states. Red signals cause queue formation.
    \item Vehicles may choose:
    \begin{itemize}
        \item a fixed path at start,
        \item or a dynamic path that changes if streets become congested (optional).
    \end{itemize}
    \item Automated/connected vehicles may be included to improve efficiency (optional).
\end{itemize}

\section*{Part 3: Outputs (Python)}

During the simulation, the program must produce:
\begin{itemize}
    \item Total wait time of all vehicles:
    \begin{itemize}
        \item at intersections (red light delays),
        \item inside edges due to crowded blocks.
    \end{itemize}
    \item Trajectory plots for selected lanes.
    \item Queue length plots for intersections over simulation time.
    \item Summary statistics (mean delays, travel times, throughput).
\end{itemize}

\section*{Part 4: SUMO Implementation}

After completing the Python simulation:
\begin{itemize}
    \item The same network (nodes and edges) must be exported or recreated in SUMO.
    \item The same demand file must be used for vehicle input.
    \item Traffic lights should follow the same logic as Python implementation.
    \item SUMO output will be collected in real-time (more realistic dynamics).
    \item Final comparison between:
    \begin{itemize}
        \item Python simulation,
        \item SUMO simulation.
    \end{itemize}
\end{itemize}

\section*{Final Deliverables}

\begin{enumerate}
    \item Python implementation with clear OOP structure (vehicles, network components, traffic lights, statistics).
    \item SUMO network configuration (\texttt{.net.xml}, \texttt{rou.xml}, traffic-light definitions).
    \item Plots: delays, queues, trajectories.
    \item A final mini-report including:
    \begin{itemize}
        \item methodology,
        \item code explanation,
        \item results from Python,
        \item results from SUMO,
        \item comparison and conclusion.
    \end{itemize}
\end{enumerate}

\section*{Grading}

\begin{itemize}
    \item Correctness of simulation logic,
    \item Object-oriented design quality,
    \item Ability to handle large demand without errors,
    \item Quality of plots and interpretation,
    \item SUMO results and comparison.
\end{itemize}

\end{document}
